\section{Completion of the reductions}\label{sec:conclusion}

\begin{proof}[Proof of Theorem~\ref{main:generic}]
Proposition~\ref{gen:genericity} gives $T=V$.
Proposition~\ref{cyc:equivalence} proves the equivalence of
conditions \textup{(i)} and \textup{(ii)}.
From condition \textup{(i)}, Theorem~\ref{gen:generation}
realizes the full primitive Clifford multiplication.
Proposition~\ref{data:polarization} and
Lemma~\ref{cliff:slots} then give condition \textup{(iii)}
with the stipulated data.
Conversely, condition \textup{(iii)} gives condition \textup{(i)}
by taking $B=A_S\times A_S$ and $\alpha=\Gamma_S$.
\end{proof}

\begin{proof}[Proof of Theorem~\ref{main:specialization}]
By Proposition~\ref{cyc:equivalence}, either assumed input gives
the nonzero correspondence very generally.
Corollary~\ref{spec:primitive}, followed by the inverse right
multiplier in Lemma~\ref{cliff:slots}, algebraically realizes
$m^q$ on $V$ at the desired point.
Proposition~\ref{data:restriction} gives $m^{q|_T}$ on $T$.
Proposition~\ref{data:polarization} and
Lemma~\ref{cliff:slots} then give the prescribed quadratic form,
parity identification, duality, and isogenous realization,
hence \eqref{eq:cycle-target}.
The operations on the abelian factors preserve the degree shift,
so the resulting cycle has codimension two.
\end{proof}

The conclusions include Picard jumps, extra real or complex multiplication, and the smallest possible transcendental ranks. They require the prescribed multiplication tensor at those points; they do not assert generation of every Hodge embedding for a smaller special Hodge group. The mathematical content of the reduction is conditional: producing a very-general nonzero algebraic correspondence, or equivalently the surface cover of Theorem~\ref{main:generic}, remains the external existence problem.
